\section{Excitons and core states}\label{app:E.15}

In silicon an electron and a hole attract each other strongly enough to bind \SI{14.7}{meV} below the gap~\cite{shaklee1970valley}, and the same attraction enhances the continuum just above it. Every near-threshold rate that a random-phase calculation gives misses both effects. The finite-momentum Bethe--Salpeter equation (BSE) of \cref{sec:4.9,app:E.6} contains them, and \cref{sec:10.5} shows how its output enters a rate. Published dark-matter calculations use its Tamm--Dancoff form, however, and none treats the bound lines of an indirect gap. We need the BSE with its coupling block, and with it the exciton series, the phonon replicas of \cref{sec:13.8}, the exciton lines of the Migdal effect (\cref{sec:12.3}) and the self-trapped exciton of an alkali halide; the semicore response and the spin--orbit mixing at the band edges (\cref{sec:10.6}) complete the set.

\subsection{Finite-momentum BSE with coupling}

What a rate near the gap needs is the loss function $L(\vect{q},\omega)$ of \cref{sec:10.2} on bins of $(|\vect{q}|,\hat{\vect{q}},\omega)$ for $|\vect{q}| \le \SI{10}{keV}$ and $\omega \le \Eg + \SI{10}{eV}$, with the energies and weights of the bound lines listed separately. Outside that box the RPA with local fields of \cref{app:E.16} serves, and the two are matched on a strip where both are computed; their difference is the many-body nuisance of \cref{sec:8.8}. On the truncation label $\mathfrak{T}$ (\cref{sec:8.5}) the calculation sits at $n_B = 0$ and $n_g = 0$, with $\mathcal{X}$ = BSE on quasiparticle bands.

At finite $q$ the coupling block between pairs and antipairs carries the long-range exchange $v_{\mathrm{C}}(\vect{q}+\vect{G})$, which turns the proper response into the screened loss function (\cref{sec:4.9}). The largest error comes from a resonant-only response, which keeps one time ordering: in silicon at $|\vect{q}| = \SI{1.39}{keV}$ and $\omega = \SI{1.3}{eV}$, in the momentum range of the phonon-assisted window, it overstates the loss function by a factor of 2.1 (\cref{eq:10.5-res}). A retarded Tamm--Dancoff calculation has no such error, since it keeps both poles $\pm E_\lambda$ and takes $\varepsilon_1$ by Kramers--Kronig; what it drops is the coupling block itself, an error set by the ratio of the local-field exchange to the pair energy, about \SI{10}{\percent} on the proper $\varepsilon_2$ (\cref{sec:10.5}). The error is avoided by the full coupling, or by keeping the Tamm--Dancoff BSE for the proper $\varepsilon_2$ (line shapes and bound lines) and taking the external screening from the RPA with local fields, with the head removed from the BSE exchange so that it is not counted twice.

The bands come from Quantum ESPRESSO with PBE and $G_0W_0$ quasiparticle energies (\cref{app:E.3}), and a code with a finite-momentum BSE, such as BerkeleyGW~\cite{deslippe2012berkeleygw}, solves the pair problem on them. For germanium the PBE starting point has $\mathrm{\Gamma}$ below L and a nearly closed gap, so the quasiparticle bands must first pass the check of \cref{app:E.14} (valley order, both gaps, masses at L and $\mathrm{\Gamma}$); the germanium BSE, and the exciton dispersion along $\mathrm{\Gamma}$--L below, wait for bands that pass. The pair basis has 4 valence bands (for germanium plus the 3d shell) and 12 conduction bands. On a coarse $8^3$ grid the kernel holds $8^3\times4\times12 = 24{,}576$ pairs, about \SI{10}{GB} per reduced momentum, and it is interpolated to a fine $24^3$ grid of $6.6\times10^5$ pairs, whose dense kernel would hold $4.4\times10^{11}$ complex numbers, \SI{7}{TB} per momentum. That kernel is never stored: a Haydock (Lanczos) solver of about 200 iterations applies it to a vector on the fly. The momenta are the reduced points of the $24^3$ set with reciprocal-lattice shells up to \SI{10}{keV}, about $10^3$ points after the cubic symmetry, each of order $10^{13}$ to $10^{14}$ floating-point operations, so the whole run takes about $10^3$ node-hours, a few times $10^4$ core-hours.

Even the fine grid cannot resolve the silicon exciton. Its radius is $a_x = \SI{4.2}{nm}$ (\cref{sec:10.5}), and a hydrogenic envelope needs a $\vect{k}$ spacing below about $1/(3a_x) \approx \SI{0.008}{\angstrom^{-1}}$, against \SI{0.08}{\angstrom^{-1}} on the $24^3$ grid. The bound lines are therefore added from the effective-mass model of \cref{sec:10.5}, with ab initio masses, $\varepsilon_\infty = 11.7$ and the vertices and their gradients at the band extrema, and the BSE supplies the continuum and its Sommerfeld enhancement. The pair basis is evaluated around the off-axis pair optimum of \cref{sec:13.9}, where the hole sits off $\mathrm{\Gamma}$, not only at the band extrema.

The coupling block restores the $f$-sum rule of \cref{sec:4.3}, which the scissor and the band cutoff also strain, so with the coupling on, the sum rule of the full loss function holds to better than \SI{1}{\percent}. At $q < 3\,\alpha m_e$ the loss function is compared with inelastic x-ray and electron energy-loss data for silicon, and the bound-line model has to reproduce the free-exciton binding of silicon, \SI{14.7}{meV}. The structure factor is reported three ways, with the full coupling, in the retarded Tamm--Dancoff form and in the resonant-only form, at $|\vect{q}| = 0.5$, 0.7, 1.0 and \SI{1.4}{keV}; the Tamm--Dancoff route passes if it agrees with the full one to \SI{15}{\percent}, and the resonant-only result tests \cref{eq:10.5-res}.

\subsection{Exciton dispersion and the hydrogenic series}

A bound exciton of total momentum $\vect{Q}$ sits below the lowest pair at that momentum by its binding energy,
\begin{equation}
  E_x(\vect{Q}) = E_{\mathrm{min}}(\vect{Q}) - E_x^b[\mu(\vect{Q}),\varepsilon_\infty] ,
  \label{eq:E.15-disp}
\end{equation}
where the reduced-mass tensor $\mu(\vect{Q})^{-1} = \nabla^2\varepsilon_c(\vect{k}_v+\vect{Q}) - \nabla^2\varepsilon_v(\vect{k}_v)$ is evaluated at the crystal momenta $\vect{k}_v(\vect{Q})$ and $\vect{k}_v(\vect{Q})+\vect{Q}$ of the lowest pair. Since the hole is not at $\mathrm{\Gamma}$ near the zero-phonon optimum (\cref{sec:13.9}), the mass changes along the dispersion. The Hessians of the Wannier-interpolated bands are taken by finite differences with a step of $0.005\times2\pi/a$, on the momentum grid of the zero-phonon edge along $\mathrm{\Gamma}$--X on both sides of X and across the $\{001\}$ planes (along $\mathrm{\Gamma}$--L and near $\mathrm{\Gamma}$ for germanium, on bands that pass the check of \cref{app:E.14}). A variational anisotropic hydrogen problem in a basis of 30 anisotropic Gaussians, with the measured $\varepsilon_\infty = 11.7$, then gives $E_x(\vect{Q})$ with the weights of the 1s and 2p states, converged when $E_x^b$ is stable to \SI{0.1}{meV} in the basis size. Near $\mathrm{\Gamma}$ the heavy and light holes are degenerate, so scalar masses fail and the hole needs the six-band $\vect{k}\cdot\vect{p}$ form with spin--orbit coupling (split-off energy \SI{44}{meV} in silicon). The calculation takes minutes; it enters \cref{sec:10.5}, and its bound lines feed the BSE above.

\subsection{Exciton--phonon replicas}

The replicas of \cref{sec:13.8}, lines at $E_x + \omega_\nu(\vect{Q}^*)$ below the free-pair edge, need the exciton--phonon vertex: the electron and hole vertices $g$ weighted by the envelopes of the initial and final excitons. Their strength per phonon branch and their position are wanted at pair momenta near $0.85\,\mathrm{X}$ and, following \cref{sec:13.9}, at the off-axis momenta where the hole sits off $\mathrm{\Gamma}$. A finite-momentum BSE at that pocket set gives them, with the vertices $g$ from EPW in the same Wannier gauge $U(\vect{k})$ as the exciton eigenvectors (\cref{app:E.14}); about $10^3$ $\vect{k}$ points per pocket resolve an exciton of radius \SI{4.2}{nm}, at about $10^4$ core-hours. A cheaper route, a Wannier--Mott envelope multiplying the free-pair amplitude of \cref{app:E.14}, is the cross-check, and both must reproduce the replica edge at \SI{0.3332}{MeV} (\cref{sec:13.9}).

Whether the replica formula of \cref{sec:13.8} is complete cannot be tested in a crystal. The band sum of \cref{sec:13.2} has one pair in every intermediate state; states with two pairs enter only through the interaction beyond the Tamm--Dancoff space of one pair (\cref{sec:13.8}), and whether the two-pair sector reduces to the one-pair terms when the electron and hole interact is proven only without interaction. A model that can be solved exactly decides it. Take a one-dimensional two-band chain of 8 to 12 cells, with spinless electrons in two cosine bands of hopping $t$ separated by $\Eg = 2t$. Add a contact attraction $W$ between electron and hole, and one phonon mode of frequency $0.05\,\Eg$ coupled linearly with strength $g$, truncated at three phonons (checked at four). The Lanczos method diagonalizes it in the sectors with zero, one and two pairs. A linear fit in $g$ between $10^{-3}$ and $10^{-2}$ extracts the exact first-order phonon-assisted amplitude, which is compared with the replica formula built from the model's exact one-pair excitons. At $W = 0$ the two must agree to $10^{-6}$, as the theorem requires. For $W/\Eg = 0.05$ to 0.3 the deviation is the two-pair correction, to be compared with its estimated size $\sim E_x^b/\Eg$, about \SI{1}{\percent} in silicon; it is extrapolated in the chain length (8, 10, 12) with a Lanczos residual below $10^{-10}$. The two-pair sector must be in the basis, whose dimension grows as $N^4$, since without it the test is circular. It runs in minutes to an hour on one core.

\subsection{Exciton lines in the Migdal response}

A nuclear kick can leave the crystal in a bound exciton rather than an ionized pair (\cref{sec:12.3}). The Migdal convolution of \cref{app:E.18} integrates the electronic momentum $\vect{p}$ against the lattice response at $\vect{q} - \vect{p}$ over the whole zone, so the exciton lines are needed everywhere, not only near $\vect{p}\to0$: the line strengths $\rho_{\lambda0}(\vect{p}+\vect{G})$, including the cross terms between different $\vect{G}$ and $\vect{G}'$, and the energies $E_\lambda(\vect{p})$. The boost operator carries the unreduced $\vect{p}$, so the Umklapp copies $\vect{p}+\vect{G}$ carry the same exciton band with different strengths, and a reduced $\vect{p}$ would merge them. The finite-momentum BSE above, in the resonant form for the proper lines with external screening from the RPA with local fields, supplies them on the 29 irreducible points of an $8^3$ grid of $\vect{p}$ and their Umklapp copies to $|\vect{p}| \le \SI{5}{keV}$, about 60 points, with the three lowest exciton bands per point, up to 20 lines and 50 reciprocal-lattice vectors. The line strengths are interpolated in $\vect{p}$ and convolved with the lattice response of \cref{app:E.17}. They must reproduce the optical exciton strength as $\vect{p}\to0$, close the $f$-sum rule with the continuum to \SI{1}{\percent} and approach the Wannier--Mott line weights in that limit, and change by less than \SI{5}{\percent} when the $\vect{k}$ points per pocket go from $10^3$ to $8\times10^3$ and when the bands are doubled. At about one node-hour per momentum the lines cost 60 to 100 node-hours, or 2 to $3\times10^3$ core-hours.

\subsection{Self-trapped excitons}

In an alkali halide such as NaI the hole localizes on a pair of anions, the lattice relaxes around it, and the energy of the self-trapped state differs from that of the free pair. Under linear coupling, a harmonic lattice and the Condon approximation, dark matter creates the exciton $\lambda$ vertically, and its line is dressed by a Poisson series of phonons~\cite{huang1950theory},
\begin{equation}
  S_\lambda(\vect{q},\omega) = |\rho_{\lambda0}(\vect{q})|^2\sum_{n\ge0}\ee^{-S_{\mathrm{HR}}}\frac{S_{\mathrm{HR}}^n}{n!}\,2\pi\delta(\omega - E_\lambda - n\bar\omega) ,
  \label{eq:E.15-ste}
\end{equation}
with the Huang--Rhys factor $S_{\mathrm{HR}}$ and a mean phonon energy $\bar\omega$; for $S_{\mathrm{HR}} \gg 1$ it is a Gaussian of mean $E_\lambda + S_{\mathrm{HR}}\bar\omega$ and variance $S_{\mathrm{HR}}\bar\omega^2$. The relaxation into the self-trapped state comes after the creation and belongs to the cascade of \cref{sec:18.5}, not to the vertex.

A finite-$q$ BSE for NaI, with the Na 2p and I 4d semicore frozen and a $6^3$ grid interpolated to $18^3$, gives $\rho_{\lambda0}(\vect{q})$ and $E_\lambda$. The relaxation energy comes from a constrained-occupation hybrid calculation of the self-trapped exciton (a hole on an anion pair plus an electron) in 64- to 216-atom supercells at the measured lattice constant, with exact-exchange fractions 0.25 and 0.40, $\mathrm{\Gamma}$-point sampling checked on a $2^3$ grid in the 64-atom cell, and the cutoff of the SG15 sodium and iodine potentials converged to \SI{1}{meV} per atom; $S_{\mathrm{HR}}$ is that energy divided by $\bar\omega$, taken from the projected phonon density of states. A semilocal functional delocalizes the hole and must not be used, and the dependence on the exchange fraction is reported as the nuisance. Iodine's spin--orbit coupling, about \SI{0.9}{eV}, splits the exciton and is included. The result is trusted when $S_{\mathrm{HR}}$ and $\bar\omega$ change by less than \SI{10}{\percent} from 64 to 216 atoms (otherwise 512 are added), when constrained occupations and time-dependent DFT give relaxation energies within \SI{20}{\percent} of each other, and when the computed luminescence energy of the self-trapped exciton matches the measured one to \SI{0.2}{eV}. The hybrid relaxations of the 216-atom cell with iodine spin--orbit coupling dominate the cost, about $10^4$ core-hours, and the time-dependent check doubles it; the factors enter \cref{sec:10.5}.

\subsection{Core levels with spin--orbit coupling}

The semicore shells of germanium and gallium, 3d at about \SI{30}{eV} below the valence maximum, and the silicon 2p shell contribute to $L(\vect{q},\omega)$ between 30 and \SI{200}{eV}, and relativistic operators enter at order $\alpha^2$ (\cref{sec:10.6}). The core part of the loss function is wanted for $q \le \SI{30}{keV}$ and $\omega \le \SI{200}{eV}$, resolved into the $3d_{5/2}$ and $3d_{3/2}$ components. An all-electron calculation with spin--orbit coupling in the zeroth-order regular approximation, in a full-potential code such as WIEN2k or exciting, or PAW with core reconstruction (\cref{app:E.2}), followed by a core-level BSE at finite $q$ on a $12^3$ grid, provides it. No crystal code computes the four-component response with projection onto positive-energy states; the two-component calculation misses only terms of order $(Z_{\mathrm{eff}}\alpha)^4 \approx 10^{-6}$ for the germanium 3d shell and a projection error that is at most \SI{2}{\percent} there, with a central estimate near $10^{-3}$ (\cref{sec:10.6}). The binding energy of the 3d shell is computed with PBE, HSE06, $G_0W_0$ and total-energy differences, and their spread about the measured doublet, $3d_{5/2}$ at \SI{29.2}{eV} and $3d_{3/2}$ at \SI{29.8}{eV} with the centroid at \SI{29.4}{eV} (\cref{sec:10.6}), is the nuisance. The decisive check is the mass--velocity shift of the $f$-sum rule of germanium, $-0.58\,\%$, the mass--velocity part of the total relativistic shift of $-0.76\,\%$ in \cref{sec:10.6}: this crystal calculation and the atomic one of \cref{app:E.16} both compute it, and they have to agree to \SI{0.1}{\percent}. The step from $12^3$ to $16^3$ must change the shift by less than \SI{0.05}{\percent} and $S(q,\omega)$ at 10 and \SI{30}{keV} by less than \SI{2}{\percent}. By analogy with the valence BSE, each material costs about $10^3$ core-hours.

\subsection{Spin--orbit mixing at the band edges}

Spin--orbit coupling mixes spin into the band edges and can lift the symmetry zeros of the density vertex. In silicon the on-axis zero is lifted to $|\rho|^2 \approx 3.6\times10^{-6}$ (summed over the four $\mathrm{\Gamma}_8$ states) and the mirror-plane zero to $1.3\times10^{-4}$ (\cref{sec:13.9}), too small to move the window. For germanium the mixing near the L-valley edge is estimated at up to \SI{10}{\percent} in amplitude, and it needs bands in the right order, which neither PBE with spin--orbit coupling nor HSE06 with it gives (\cref{app:E.14}). The step therefore runs on germanium bands with spin--orbit coupling that pass that check. On those bands a non-self-consistent run at the crystal momenta $(\vect{k}_v, \vect{k}_v+\vect{q})$ of the zero-phonon optimum gives the vertices from spinor wavefunctions matched by Miller index. Single spinor bands are gauge dependent, so only sums over Kramers pairs, or over degenerate quartets, are compared. Those sums are converged when they change by less than \SI{2}{\percent} from 60 to \SI{80}{Ry} and when the L and $\mathrm{\Gamma}$ energies of the underlying bands are stable to \SI{10}{meV} in the number of bands and the screening cutoff. The silicon result, which must vanish as the spin--orbit strength is scaled to zero, is the control. Once bands that pass exist, the step costs less than a core-hour, and it replaces the estimate of \cref{sec:10.6}.
